\clearpage
\SourceLeaf{179}
\pdfbookmark[0]{Göttingische gelehrte Anzeigen. 1831, 23. April.}{notice1831april23}
\begin{center}
\rule{.82\linewidth}{.35pt}\par
Göttingische gelehrte Anzeigen.\quad 1831, 23. April.\par
\rule{.82\linewidth}{.35pt}
\end{center}
Eine Vorlesung, die Hofrat \textsc{Gauss} der Königlichen Gesellschaft am 15. April überreicht hat:
\begin{center}\itshape
Theoria residuorum biquadraticorum, commentatio secunda,
\end{center}
setzt die bereits im sechsten Band der \emph{Commentationes novae} abgedruckte Abhandlung fort, die seinerzeit auch in unseren Blättern, im Jahr 1825, am 11. April, angezeigt wurde. Auch diese Fortsetzung, obwohl mehr als doppelt so umfangreich wie die erste Abhandlung, erschöpft den überaus reichen Gegenstand noch nicht. Die Vollendung des Ganzen wird erst einer künftigen dritten Abhandlung vorbehalten bleiben.

Obwohl wir die Grundbegriffe dieser Lehren und den Inhalt der ersten Abhandlung bei allen, die höhere Arithmetik studiert haben, als bekannt voraussetzen können, wollen wir jene Begriffe hier doch kurz in Erinnerung bringen. Dies soll den Freunden dieses Teils der Mathematik dienen, die die erste Abhandlung nicht gleich zur Hand haben. In Bezug auf eine beliebige ganze Zahl $p$ heißt eine andere ganze Zahl $k$ ein biquadratischer Rest, wenn es Zahlen der Form $x^4-k$ gibt, die durch $p$ teilbar sind; im entgegengesetzten Fall heißt sie biquadratischer Nichtrest von $p$. Es genügt hier, sich auf den Fall zu beschränken, dass $p$ eine Primzahl der Form $4n+1$ und $k$ durch
\SourceLeaf{180}
diese Primzahl nicht teilbar ist, da alle anderen Fälle entweder für sich klar sind oder sich auf diesen zurückführen lassen.

Für einen solchen \emph{gegebenen} Wert von $p$ zerfallen sämtliche durch $p$ nicht teilbaren Zahlen in \emph{vier} Klassen. Eine enthält die biquadratischen Reste, eine zweite diejenigen biquadratischen Nichtreste, die quadratische Reste von $p$ sind; auf die beiden übrigen verteilen sich die biquadratischen Nichtreste, die zugleich quadratische Nichtreste sind. Das Prinzip dieser Aufteilung besteht darin, dass stets entweder $k^n-1$ oder $k^n+1$ oder $k^n-f$ oder $k^n+f$ durch $p$ teilbar ist. Dabei bedeutet $f$ eine ganze Zahl, die $ff+1$ durch $p$ teilbar macht. Jeder, der die elementare Terminologie kennt, sieht von selbst, wie sich diese Worterklärungen darin ausdrücken lassen.

Die Theorie dieser Klassifikation ist in der ersten Abhandlung vollständig ausgeführt: sowohl für den unmittelbar einleuchtenden Fall $k=-1$ als auch für die Fälle $k=\pm2$, die subtile Hilfsuntersuchungen erfordern. Am Anfang der gegenwärtigen Abhandlung wird nun zu größeren Werten von $k$ übergegangen. Zunächst braucht man aber nur solche Werte zu betrachten, die selbst Primzahlen sind. Dabei zeigt sich, dass die Ergebnisse am einfachsten ausfallen, wenn man die Werte positiv oder negativ nimmt, je nachdem sie ihrem Betrag nach die Form $4m+1$ oder $4m+3$ haben. Die Untersuchung von Zahlenbeispielen liefert hier sofort und mit großer Leichtigkeit eine reiche Ernte neuer Lehrsätze, von denen wir nur einige anführen. Die Nummerierung der Klassen mit $1,2,3,4$ bezieht sich auf die Fälle, in denen $k^n$ jeweils zu $1,f,-1,-f$ kongruent ist. Zugleich wird für $f$ stets derjenige Wert angenommen, der $a+bf$ durch $p$ teilbar macht, wenn $aa+bb$ die Zerlegung von $p$ in ein ungerades und ein gerades Quadrat darstellt. So findet man durch die Untersuchung von Zahlenbeispielen, dass $-3$ jeweils zur Klasse $1,2,3,4$ gehört, je nachdem $b,a+b,a,a-b$ durch $3$ teilbar ist; dass $+5$ der Reihe nach zu diesen Klassen gehört, je nachdem $b,a-b,a,a+b$ durch $5$ teilbar ist; dass $-7$ in die Klasse $1$ fällt, wenn $a$ oder $b$; in die Klasse $2$, wenn $a-2b$ oder $a-3b$; in die Klasse $3$, wenn $a-b$ oder $a+b$; in die Klasse $4$, wenn $a+2b$ oder $a+3b$ durch $7$ teilbar ist, wobei die Teilbarkeit in jedem dieser Fälle die Bedingung ist. Ähnliche Sätze ergeben sich für die Zahlen $-11,+13,+17,-19,-23$ usw. So leicht sich aber all diese besonderen Sätze durch die Untersuchung von Zahlenbeispielen entdecken lassen, so schwer scheint
\SourceLeaf{181}
es, auf diesem Weg ein allgemeines Gesetz für diese Formen zu finden, auch wenn manches Gemeinsame bald auffällt. Noch viel schwerer ist es, Beweise für diese Lehrsätze zu finden. Die in der ersten Abhandlung für die Zahlen $+2$ und $-2$ verwendeten Methoden lassen sich hier nicht mehr anwenden. Andere Methoden könnten zwar ebenfalls das erledigen, was die erste und dritte Klasse betrifft, erweisen sich aber zur Begründung \emph{vollständiger} Beweise als ungeeignet.

Man erkennt daher bald, dass dieses reiche Gebiet der höheren Arithmetik nur auf ganz neuen Wegen zugänglich ist. Der Verfasser hatte schon in der ersten Abhandlung angedeutet, dass dazu eine eigentümliche Erweiterung des gesamten Feldes der höheren Arithmetik wesentlich erforderlich ist, ohne damals näher zu erklären, worin sie besteht. Die gegenwärtige Abhandlung soll diesen Gegenstand ins Licht setzen.

Es geht um nichts anderes, als dass für die wahre Begründung der Theorie der biquadratischen Reste das Feld der höheren Arithmetik, das man sonst nur auf die reellen ganzen Zahlen erstreckte, auch auf die imaginären ausgedehnt werden muss. Diesen muss das völlig gleiche Bürgerrecht wie jenen eingeräumt werden. Sobald man das eingesehen hat, erscheint jene Theorie in einem ganz neuen Licht, und ihre Ergebnisse gewinnen eine höchst überraschende Einfachheit.

Bevor jedoch in diesem erweiterten Zahlengebiet die Theorie der biquadratischen Reste selbst entwickelt werden kann, müssen auch die ihr vorangehenden Lehren der höheren Arithmetik an dieser Erweiterung teilnehmen, die bisher nur für reelle Zahlen behandelt worden sind. Von diesen vorbereitenden Untersuchungen können wir hier nur einiges anführen. Der Verfasser nennt eine Größe $a+bi$, in der $a$ und $b$ reelle Größen bedeuten und $i$ der Kürze wegen für $\surd-1$ steht, eine komplexe ganze Zahl, wenn zugleich $a$ und $b$ ganze Zahlen sind. Die komplexen Größen stehen also nicht den reellen gegenüber, sondern enthalten diese als den besonderen Fall, dass $b=0$. Für einen bequemeren Gebrauch mussten mehrere Begriffsbildungen zu komplexen Größen eigene Benennungen erhalten, die wir in dieser Anzeige jedoch zu umgehen versuchen werden.

So wie in der Arithmetik der reellen Zahlen nur von zwei Einheiten, der positiven und der negativen, die Rede ist, so haben wir in der Arithmetik der
\SourceLeaf{182}
komplexen Zahlen vier Einheiten: $+1$, $-1$, $+i$, $-i$. Eine komplexe ganze Zahl heißt \emph{zusammengesetzt}, wenn sie das Produkt zweier ganzzahliger Faktoren ist, die keine Einheiten sind. Eine komplexe Zahl, die eine \emph{solche} Zerlegung in Faktoren nicht zulässt, heißt dagegen eine komplexe Primzahl. So ist beispielsweise die reelle Zahl $3$ auch als komplexe Zahl eine Primzahl, während $5$ als komplexe Zahl zusammengesetzt ist $=(1+2i)(1-2i)$. Wie in der höheren Arithmetik der reellen Zahlen spielen die Primzahlen auch in diesem erweiterten Gebiet der Wissenschaft eine Hauptrolle.

Nimmt man eine komplexe ganze Zahl $a+bi$ als Modul an, so lassen sich $aa+bb$ paarweise nicht kongruente komplexe Zahlen angeben, und nicht mehr, zu deren einer jede vorgegebene ganze komplexe Zahl kongruent sein muss. Man kann sie ein vollständiges System inkongruenter Reste nennen. Die sogenannten kleinsten und absolut kleinsten Reste in der Arithmetik der reellen Zahlen haben auch hier ihr genaues Gegenstück. So besteht beispielsweise für den Modul $1+2i$ das vollständige System der absolut kleinsten Reste aus den Zahlen $0$, $1$, $i$, $-1$ und $-i$. Fast sämtliche Untersuchungen der ersten vier Abschnitte der \emph{Disquisitiones Arithmeticae} finden mit einigen Abwandlungen auch in der erweiterten Arithmetik ihren Platz. Das berühmte Theorem von \textsc{Fermat} nimmt hier beispielsweise folgende Gestalt an: Ist $a+bi$ eine komplexe Primzahl und $k$ eine durch sie nicht teilbare komplexe Zahl, so gilt stets $k^{aa+bb-1}\equiv1$ für den Modul $a+bi$. Besonders bemerkenswert ist aber, dass das Fundamentaltheorem für die quadratischen Reste in der Arithmetik der komplexen Zahlen sein genaues und hier sogar einfacheres Gegenstück hat. Sind nämlich $a+bi$, $A+Bi$ komplexe Primzahlen, wobei $a$ und $A$ ungerade, $b$ und $B$ gerade sind, so ist die erste quadratischer Rest der zweiten, wenn die zweite quadratischer Rest der ersten ist; dagegen ist die erste quadratischer Nichtrest der zweiten, wenn die zweite quadratischer Nichtrest der ersten ist.

Wenn die Abhandlung nach diesen Voruntersuchungen zur Lehre von den biquadratischen Resten selbst übergeht, legt sie zunächst anstelle der bloßen Unterscheidung zwischen biquadratischen Resten und Nichtresten eine Aufteilung der durch den Modul nicht teilbaren Zahlen in vier Klassen fest. Ist nämlich der Modul eine komplexe Primzahl $a+bi$, wobei stets $a$ ungerade und $b$ gerade vorausgesetzt wird und der Kürze wegen $p$ für $aa+bb$ steht, und ist $k$ eine komplexe, durch $a+bi$ nicht teilbare Zahl, so ist stets $k^{\frac14(p-1)}$
\SourceLeaf{183}
zu einer der Zahlen $+1,+i,-1,-i$ kongruent. Dadurch werden sämtliche durch $a+bi$ nicht teilbaren Zahlen in vier Klassen aufgeteilt, denen der Reihe nach der biquadratische Charakter $0,1,2,3$ zugeordnet wird. Offenbar bezieht sich der Charakter $0$ auf die biquadratischen Reste, die übrigen auf die biquadratischen Nichtreste, und zwar so, dass dem Charakter $2$ zugleich quadratische\textperiodcentered\ Reste entsprechen, den Charakteren $1$ und $3$ dagegen quadratische Nichtreste.

Man erkennt leicht, dass es vor allem darauf ankommt, diesen Charakter nur für solche Werte von $k$ bestimmen zu können, die selbst komplexe Primzahlen sind. Hier führt die Untersuchung von Zahlenbeispielen sofort zu höchst einfachen Ergebnissen.

Setzt man zunächst $k=1+i$, so zeigt sich, dass der Charakter dieser Zahl stets $\equiv\frac18(-aa+2ab-3bb+1)(\mathrm{mod.}\ 4)$ wird. Ähnliche Ausdrücke findet man für die Fälle $k=1-i$, $k=-1+i$, $k=-1-i$.

Ist dagegen $k=\alpha+\beta i$ eine Primzahl mit ungeradem $\alpha$ und geradem $\beta$, so ergibt sich durch die Untersuchung von Zahlenbeispielen sehr leicht ein zum Fundamentaltheorem für die quadratischen Reste ganz analoges Reziprozitätsgesetz. Es lässt sich am einfachsten folgendermaßen ausdrücken:

Sind sowohl $\alpha+\beta-1$ als auch $a+b-1$ durch $4$ teilbar (auf diesen Fall lassen sich alle übrigen leicht zurückführen), und bezeichnet man den Charakter von $\alpha+\beta i$ in Bezug auf den Modul $a+bi$ mit $\lambda$, dagegen den Charakter von $a+bi$ in Bezug auf den Modul $\alpha+\beta i$ mit $l$, so gilt $\lambda=l$, wenn eine der Zahlen $\beta,b$ (oder beide) durch $4$ teilbar ist. dagegen $\lambda=l\pm2$, wenn keine der Zahlen $\beta,b$ durch $4$ teilbar ist.

Diese Sätze enthalten im Grunde alles Wesentliche der Theorie der biquadratischen Reste. So leicht es aber war, sie durch die Untersuchung von Zahlenbeispielen zu entdecken, so schwer ist es, strenge Beweise für sie zu geben, besonders für den zweiten, das Fundamentaltheorem der biquadratischen Reste. Wegen des großen Umfangs, den die gegenwärtige Abhandlung bereits erreicht hat, sah sich der Verfasser gezwungen, die Darstellung des Beweises für den letzteren Satz, den er seit $20$ Jahren besitzt, einer künftigen dritten Abhandlung vorzubehalten. Die vorliegende Abhandlung enthält dagegen noch den vollständigen Beweis des ersteren Satzes über $1+i$ (von dem die
\SourceLeaf{184}
anderen Sätze für $1-i$, $-1+i$, $-1-i$ abhängen). Dieser kann bereits eine Vorstellung von der Verwicklung des Gegenstandes geben.

Wir haben noch einige allgemeine Anmerkungen hinzuzufügen. Die Verlegung der Lehre von den biquadratischen Resten in das Gebiet der komplexen Zahlen könnte manchem, der mit der Natur der imaginären Größen weniger vertraut und in falschen Vorstellungen über sie befangen ist, anstößig und unnatürlich erscheinen. Sie könnte die Meinung hervorrufen, die Untersuchung werde dadurch gleichsam in die Luft gestellt, erhalte einen schwankenden Halt und entferne sich ganz von der Anschaulichkeit. Nichts wäre unbegründeter als eine solche Meinung. Im Gegenteil lässt die Arithmetik der komplexen Zahlen die anschaulichste Veranschaulichung zu. Der Verfasser hat in seiner diesmaligen Darstellung zwar eine rein arithmetische Behandlung gewählt, aber auch die nötigen Hinweise zu dieser Veranschaulichung gegeben, die die Einsicht lebendiger macht und deshalb sehr zu empfehlen ist. sie werden selbstständig denkenden Lesern genügen. So wie die absoluten ganzen Zahlen durch eine Reihe gleich weit voneinander entfernter Punkte auf einer Geraden dargestellt werden, in der der Anfangspunkt die Zahl $0$, der nächste die Zahl $1$ usw. vertritt; und so wie man zur Darstellung der negativen Zahlen diese Reihe nur auf der entgegengesetzten Seite des Anfangspunkts unbegrenzt verlängern muss: so braucht man zur Darstellung der komplexen ganzen Zahlen lediglich hinzuzufügen, dass jene Reihe in einer bestimmten unbegrenzten Ebene liegt und parallel zu ihr auf beiden Seiten eine unbeschränkte Anzahl ähnlicher Reihen in gleichen Abständen voneinander angenommen wird. Statt einer Reihe von Punkten haben wir dann ein System von Punkten vor uns, die sich auf zwei Arten in Reihen von Reihen ordnen lassen und die ganze Ebene in gleiche Quadrate einteilen. Der nächste Punkt bei $0$ in der ersten Nebenreihe auf der einen Seite der Reihe, die die reellen Zahlen darstellt, bezieht sich dann auf die Zahl $i$, ebenso der nächste Punkt bei $0$ in der ersten Nebenreihe auf der anderen Seite auf $-i$ usw. Bei dieser Darstellung lassen sich die Ausführung der arithmetischen Operationen mit komplexen Größen, die Kongruenz, die Bildung eines vollständigen Systems inkongruenter Zahlen für einen gegebenen Modul usw. in einer Weise veranschaulichen, die nichts zu wünschen übrig lässt.

\SourceLeaf{185}
Andererseits wird dadurch die wahre Metaphysik der imaginären Größen in ein neues helles Licht gestellt.

Unsere allgemeine Arithmetik, deren Umfang die Geometrie der Alten so weit übertrifft, ist ganz eine Schöpfung der neueren Zeit. Ursprünglich vom Begriff der absoluten ganzen Zahlen ausgehend, hat sie ihr Gebiet schrittweise erweitert: Zu den ganzen Zahlen kamen die gebrochenen, zu den rationalen die irrationalen, zu den positiven die negativen, zu den reellen die imaginären hinzu. Jeder dieser Fortschritte geschah aber zunächst mit furchtsam zögernden Schritten. Die ersten Algebraiker nannten die negativen Wurzeln der Gleichungen noch falsche Wurzeln. Das sind sie auch, wenn die Aufgabe, auf die sie sich beziehen, in einer Form gestellt ist, in der die gesuchte Größe ihrer Beschaffenheit nach kein Entgegengesetztes zulässt. Aber so wenig man in der \emph{allgemeinen} Arithmetik Bedenken hat, gebrochene Zahlen aufzunehmen, obwohl es so viele zählbare Dinge gibt, bei denen eine Bruchzahl sinnlos ist, ebenso wenig durfte man dort den negativen Zahlen die gleichen Rechte wie den positiven deshalb versagen, weil unzählige Dinge kein Entgegengesetztes zulassen. Die Realität der negativen Zahlen ist hinreichend gerechtfertigt, da sie in unzähligen anderen Fällen ein angemessenes Substrat finden. Darüber besteht freilich schon seit langer Zeit Klarheit. Die den reellen Größen gegenübergestellten imaginären dagegen---früher. und hier und da noch heute, wenn auch unpassend, \emph{unmögliche} genannt---sind noch immer weniger eingebürgert als bloß geduldet. Sie erscheinen daher eher als ein an sich inhaltsleeres Zeichenspiel, dem man ein denkbares Substrat unbedingt abspricht, ohne jedoch den reichen Tribut verschmähen zu wollen, den dieses Zeichenspiel schließlich zum Schatz der Verhältnisse reeller Größen beiträgt.

Der Verfasser betrachtet diesen höchst wichtigen Teil der Mathematik seit vielen Jahren aus einem anderen Blickwinkel, unter dem man den imaginären Größen ebenso gut einen Gegenstand zugrunde legen kann wie den negativen. Bisher fehlte jedoch ein Anlass, diese Sicht öffentlich und ausdrücklich darzulegen. Aufmerksame Leser werden ihre Spuren allerdings leicht in der $1799$ erschienenen Schrift über Gleichungen und in der Preisschrift über die Umbildung der Flächen wiederfinden. In der gegenwärtigen Abhandlung werden ihre Grundzüge kurz angegeben; sie bestehen in Folgendem.

Positive und negative Zahlen lassen sich nur dort anwenden, wo
\SourceLeaf{186}
das Gezählte ein Entgegengesetztes hat, das, mit ihm vereinigt gedacht, der Vernichtung gleichzusetzen ist. Genau betrachtet gilt diese Voraussetzung nur dort, wo nicht Substanzen (für sich denkbare Gegenstände), sondern Beziehungen zwischen jeweils zwei Gegenständen gezählt werden. Dabei wird vorausgesetzt, dass diese Gegenstände auf bestimmte Art in eine Reihe geordnet sind, beispielsweise $A,B,C,D\dotsfour$, und dass die Beziehung von $A$ zu $B$ als gleich der Beziehung von $B$ zu $C$ usw. gelten kann. Zum Begriff der Entgegensetzung gehört hier nichts weiter als der \emph{Austausch} der Glieder der Beziehung. Gilt also die Beziehung (oder der Übergang) von $A$ zu $B$ als $+1$, so muss die Beziehung von $B$ zu $A$ durch $-1$ dargestellt werden. Soweit eine solche Reihe auf beiden Seiten unbegrenzt ist, stellt also jede reelle ganze Zahl die Beziehung eines beliebig als Anfang gewählten Gliedes zu einem bestimmten Glied der Reihe dar.

Sind die Gegenstände aber so beschaffen, dass sie sich nicht in eine einzige, wenn auch unbegrenzte Reihe ordnen lassen, sondern nur in Reihen von Reihen, oder, gleichbedeutend damit, bilden sie eine Mannigfaltigkeit von zwei Dimensionen; und verhalten sich die Beziehungen einer Reihe zu einer anderen oder die Übergänge zwischen ihnen ähnlich wie zuvor die Übergänge von einem Glied einer Reihe zu einem anderen Glied derselben Reihe, so braucht man offenbar zum Messen des Übergangs von einem Glied des Systems zu einem anderen außer den bisherigen Einheiten $+1$ und $-1$ noch zwei weitere, einander ebenfalls entgegengesetzte Einheiten $+i$ und $-i$. Dabei muss aber noch vorausgesetzt werden, dass die Einheit $i$ stets den Übergang von einem gegebenen Glied einer Reihe zu einem \emph{bestimmten} Glied der unmittelbar angrenzenden Reihe bezeichnet. Auf diese Weise lässt sich das System also auf zwei Arten in Reihen von Reihen ordnen.

Der Mathematiker sieht von der Beschaffenheit der Gegenstände und vom Inhalt ihrer Beziehungen vollständig ab. Er hat es nur damit zu tun, die Beziehungen zu zählen und miteinander zu vergleichen. Insofern darf er die Gleichartigkeit, die er den durch $+1$ und $-1$ bezeichneten Beziehungen an sich zuschreibt, auf alle vier Elemente $+1$, $-1$, $+i$ und $-i$ ausdehnen.

Anschaulich werden diese Verhältnisse nur durch eine Darstellung
\SourceLeaf{187}
im Raum. Der einfachste Fall liegt vor, wenn kein Grund besteht, die Symbole der Gegenstände anders als quadratisch anzuordnen: Man teilt eine unbegrenzte Ebene durch zwei Systeme paralleler Linien, die sich rechtwinklig kreuzen, in Quadrate und wählt die Schnittpunkte als Symbole. Jeder solche Punkt $A$ hat hier vier Nachbarn. Bezeichnet man die Beziehung von $A$ zu einem benachbarten Punkt mit $+1$, so ist die mit $-1$ zu bezeichnende Beziehung von selbst bestimmt. Dagegen kann man für $+i$ nach Belieben eine der beiden anderen wählen, also den zu $+i$ gehörenden Punkt \emph{rechts} oder \emph{links} nehmen. Sobald man vorwärts und rückwärts \emph{in} der Ebene sowie oben und unten in Bezug auf ihre beiden Seiten einmal nach Belieben festgelegt hat, ist dieser Unterschied zwischen rechts und links \emph{in sich} völlig bestimmt. Unsere Anschauung dieses Unterschieds können wir anderen allerdings \emph{nur} vermitteln, indem wir auf tatsächlich vorhandene materielle Dinge zeigen\begingroup\renewcommand{\thefootnote}{*)}\footnote{Beide Bemerkungen hat bereits \textsc{Kant} gemacht. Man begreift aber nicht, wie dieser scharfsinnige Philosoph in der ersten einen Beweis für seine Ansicht finden zu können glaubte, dass der Raum \emph{nur} Form unserer äußeren Anschauung sei, da die zweite so klar das Gegenteil beweist und zeigt, dass der Raum unabhängig von unserer Anschauungsart eine reelle Bedeutung haben muss.}\endgroup. Auch wenn man sich über Letzteres entschieden hat, sieht man jedoch, dass es von unserer Willkür abhing, welche der beiden sich in einem Punkt kreuzenden Reihen wir als Hauptreihe und welche Richtung darin wir als auf positive Zahlen bezogen ansehen wollten. Ferner sieht man: Will man die zuvor als $+i$ behandelte Beziehung für $+1$ nehmen, muss man notwendig die zuvor durch $-1$ bezeichnete Beziehung für $+i$ nehmen. In der Sprache der Mathematiker heißt das: $+i$ ist eine mittlere Proportionalgröße zwischen $+1$ und $-1$ oder entspricht dem Zeichen $\surd-1$. Wir sagen absichtlich nicht \emph{die} mittlere Proportionalgröße, denn $-i$ hat offenbar den gleichen Anspruch. Damit ist die Möglichkeit, für $\surd-1$ eine anschauliche Bedeutung anzugeben, vollkommen gerechtfertigt. Mehr braucht es nicht, um diese Größe in das Gebiet der Gegenstände der Arithmetik aufzunehmen.

\newpage
Wir glaubten, den Freunden der Mathematik mit dieser kurzen Darstellung der Hauptpunkte einer neuen Theorie der sogenannten imaginären Größen einen Dienst zu erweisen. Wenn man diesen Gegenstand bisher aus einem falschen Blickwinkel betrachtet und dabei eine geheimnisvolle Dunkelheit
\SourceLeaf{188}
gefunden hat, so liegt das größtenteils an den wenig passenden Benennungen. Hätte man $+1$, $-1$, $\surd-1$ beispielsweise direkte, inverse und laterale Einheit genannt statt positive, negative und imaginäre (oder gar unmögliche) Einheit, so hätte von solcher Dunkelheit kaum die Rede sein können. Der Verfasser hat sich vorbehalten, den Gegenstand, der in der vorliegenden Abhandlung eigentlich nur nebenbei berührt wird, künftig vollständiger zu behandeln. Dabei wird auch die Frage beantwortet werden, warum die Beziehungen zwischen Dingen, die eine Mannigfaltigkeit von mehr als zwei Dimensionen bilden, nicht noch weitere Arten von Größen liefern können, die in der allgemeinen Arithmetik zulässig sind.
\BookRule{.28\linewidth}

